%% ============================================================
%% sections/a7-explainability.tex
%% H. R. 9510 Bill v3.0 - APPENDIX C. THE VVUQ EXPLAINABILITY STANDARD
%% (NON-OPERATIVE). Establishes the VVUQ-01 through VVUQ-04 developments as a
%% standard for explainability of the U.S. legislation-creation process. Rendered
%% visuals: Figure 9 (prompt evolution), Figure 10 (internal bill evolution),
%% Table 7 (the explainability standard mapped to the four NISTIR 8312 principles).
%% ============================================================
\clearpage
\phantomsection\label{toc:a7}
\addcontentsline{toc}{section}{Appendix C. The VVUQ Explainability Standard (Non-Operative)}
\usccrosshead{APPENDIX C. THE VVUQ EXPLAINABILITY STANDARD (NON-OPERATIVE)}

\uscprov{1}{This appendix is non-operative. It establishes the VVUQ-01 through
VVUQ-04 developments as a reusable standard for explaining how a piece of United
States legislation was created. The claim is narrow and demonstrable: every step
that produced this bill, from the first method paper to the precise amendment,
recorded the exact prompt that drove it, the full output it produced, a VVUQ gate
decision that held the output to its specification, and a commit trail that pushed
each file in real time. A reader can therefore reconstruct how the law was built,
which is the property an explainable process must have \cite{nist-ai-explainability,
nist-ai-600}.}

\uscprov{1}{The prompt record is the first half of the standard. Because each
stage keeps its prompt and its output side by side, the reasoning that turned
evidence into statutory language is inspectable rather than hidden, as Figure 9
shows: a reviewer can read the update-bill prompt that retargeted the work onto
the Federal Food, Drug, and Cosmetic Act, then read the output that resulted
\cite{vvuq05-figures-bill}.}

\begin{asciifig}
PROMPT EVOLUTION (each stage records the instruction that drove it)

  VVUQ-02   instruct -> draft -> full
     |      (method: build and prove the ten-gate suite)
     v
  VVUQ-03   instruct -> draft -> full
     |      (law: turn the recorded evidence into a bill)
     v
  VVUQ-04   instruct -> update -> draft -> full
     |      (amendment: prompt-a-update-bill retargets onto the FD&C Act)
     v
  VVUQ-05   figures -> next-steps -> draft -> full   <== THIS
            (visual perspective and the submission deliverables)

  the instruction sharpens from "build the gate" to "amend Title 21 precisely"
\end{asciifig}
\vspace{-0.7cm}
\figcaption{Figure 9. Prompt evolution across the works: the chain of recorded
prompts from VVUQ-02 through\\VVUQ-05, showing how the instruction to the agent
sharpened as the work moved from method to statute.}

\uscprov{1}{The stage record is the second half. The bill did not appear in one
pass: a legal crosswalk was assembled, the real statute was reproduced, the
amendment was scaffolded with bracketed instructions, the instructions were
executed, and the result was finalized, with this visual bill a further stage, as
Figure 10 shows. Each stage is a directory a reviewer can open, so the provenance
of every clause is traceable to the files it was synthesized from
\cite{kawchak2026vvuq04bill}.}

\begin{asciifig}
INTERNAL BILL EVOLUTION (each stage a directory feeding the next)

  VVUQ-04
  instruct-bill -> template-bill -> draft-bill -> full-bill -> final-bill
                                                               |
                                                               v
  VVUQ-05
  update-bill (figures-bill + next-steps) -> draft-bill -> full-bill
                                                           |
                                                           v
                                                      Bill v3.0 <== THIS
\end{asciifig}
\vspace{-0.7cm}
\figcaption{Figure 10. Internal bill evolution: the VVUQ-04 stages, extended by
the VVUQ-05 update-bill\\(figures-bill and next-steps) into this visual bill, each
stage a directory feeding the next.}

\uscprov{1}{Mapped to the four principles of explainable artificial intelligence,
the standard reads as in Table 7: explanation is met because every step has a
recorded prompt and output; meaningful is met because each stage carries a README
that explains the step to a human; accuracy is met because the VVUQ gate holds each
output to its specification and the results reproduce from a fixed seed; and the
knowledge-limits principle is met because the gate escalates to a human and the
disclaimers mark where the system must hand back. This is the explainability the
legislation-creation process should carry, demonstrated end to end
\cite{nist-ai-explainability, kawchak2026vvuq01}.}

{\footnotesize
\begin{xltabular}{\textwidth}{@{}L{2.3cm} Y L{3.4cm}@{}}
\hline\noalign{\vskip 0.04cm}
\textbf{Dimension} & \textbf{What makes the legislation-creation step satisfy it}
& \textbf{Recording artifact}\\
\hline\noalign{\vskip 0.04cm}
\endfirsthead
\hline\noalign{\vskip 0.04cm}
\textbf{Dimension} & \textbf{What makes the legislation-creation step satisfy it}
& \textbf{Recording artifact}\\
\hline\noalign{\vskip 0.04cm}
\endhead
\hline\noalign{\vskip 0.04cm}
\endlastfoot
Explanation & Every step states the exact instruction that drove it and the full
result & Prompt and output Markdown pair at stage\\
\hline\noalign{\vskip 0.04cm}
Meaningful & Each stage is explained to a human in terms a reviewer can follow &
The per-stage README\\
\hline\noalign{\vskip 0.04cm}
Accuracy & Each output is held to its specification and reproduces from a fixed
seed & The VVUQ gate decision surface; seed 20260525\\
\hline\noalign{\vskip 0.04cm}
Knowledge limits & The process hands back to a human on divergence and marks what
it is not & The ESCALATE rule and the standing disclaimers\\
\hline
\end{xltabular}}
\vspace{-0.7cm}
\figcaption{Table 7. The VVUQ explainability standard: each principle of
explainable artificial intelligence (NISTIR 8312), what makes a
legislation-creation step satisfy it, and the artifact in the repository that
records it.}
